\chapter{Phân Tích Mô Hình Hóa Hệ Thống \& Sơ Đồ Use Case}
\label{ch:use_case_modeling}
\textit{Vo Bao Phuc}

\section{Xác Định Các Tác Nhân (Actors) Vận Hành Hệ Thống}
Mô hình hóa hệ thống xác định các thực thể tương tác trực tiếp bao gồm:
\begin{itemize}
	\item \textbf{Sinh viên (Student):} Tác nhân chính thực hiện tương tác giao diện, làm bài kiểm tra, xem lộ trình và kết nối GitHub cá nhân.
	\item \textbf{Trí tuệ nhân tạo (AI Recommendation Engine):} Tác nhân hệ thống ngầm, chịu trách nhiệm nhận tham số đầu vào là hồ sơ sinh viên để sinh cây kỹ năng.
	\item \textbf{Robot quét dữ liệu (Market Pulse Bot):} Tác nhân tự động thực hiện thu thập thông tin tuyển dụng từ internet.
	\item \textbf{Quản trị viên (Admin):} Tác nhân quản lý hệ thống danh mục, kiểm duyệt nội dung và định cấu hình bộ khung năng lực chuẩn.
\end{itemize}

\section{Bảng Đặc Tả Chi Tiết Use Case Cốt Lõi}
Dưới đây là kịch bản vận hành chi tiết cho Use Case trọng tâm của hệ thống: \textit{Tạo lộ trình học tập cá nhân hóa}:

\begin{table}[H]
	\centering
	\small
	\caption{Đặc tả Use Case Tạo Lộ Trình Học Tập Cá Nhân Hóa}
	\label{tab:spec_usecase}
	\begin{tabular}{|p{3.5cm}|p{11cm}|}
		\hline
		\rowcolor{mainblue} \color{white}\textbf{Thành Phần Đặc Tả} & \color{white}\textbf{Nội Dung Mô Tả Chi Tiết} \\ \hline
		\textbf{Tên Use Case} & Tạo lộ trình học tập cá nhân hóa (Generate Personalized Roadmap) \\ \hline
		\textbf{Tác nhân (Actors)} & Sinh viên (Student), Hệ thống AI Engine (Supporting Actor) \\ \hline
		\textbf{Mục tiêu hành động} & Sinh viên có được biểu đồ cây chặng đường học tập được thiết kế riêng. \\ \hline
		\textbf{Tiền điều kiện} & Sinh viên đã hoàn tất đăng nhập và đã đồng bộ dữ liệu kỹ năng hiện có. \\ \hline
		\textbf{Luồng sự kiện chính} & 
		1. Sinh viên nhấn chọn chức năng ``Tạo lộ trình cá nhân hóa'' trên giao diện. \newline
		2. Hệ thống thu thập tập dữ liệu hồ sơ năng lực hiện tại của người dùng. \newline
		3. Hệ thống chuyển đổi dữ liệu, gọi API lớp Business Logic để tính khoảng cách năng lực. \newline
		4. Hệ thống AI Engine phân tích và trả lời cấu trúc dữ liệu cây kỹ năng còn thiếu. \newline
		5. Giao diện kết xuất đồ thị trực quan hiển thị chặng đường học tập cho sinh viên. \\ \hline
		\textbf{Hậu điều kiện} & Cấu trúc lộ trình học tập mới được lưu trữ an toàn vào cơ sở dữ liệu vật lý. \\ \hline
	\end{tabular}
\end{table}